\section{测试、故障注入与可观察性}

\subsection{测试金字塔必须落到硬件边界}

宿主单元测试验证地址范围、ROM 解码和串口协议解析；集成测试验证链接产物；
固定种子随机测试生成大量错误复位位移，确认审计器不会误接受；QEMU 系统测试
执行真实目标字节。它们不是重复劳动，而是分别缩小不同故障的定位范围。

\subsection{有界运行}

早期固件最终执行 \code{hlt}，后续内核则在 HLT 事件循环等待 IRQ。测试进程
逐行观察当前用例最后一个必需里程碑，到达后短暂收尾并回收 QEMU；尚未到达时
最多等待五秒。成功用例要求全部标记按序存在；失败用例要求对应失败标记存在，
并禁止越过失败阶段。异常提前退出同样视为失败。

\begin{failurebox}
仅检查进程“还活着”不能证明 CPU 执行了我们的代码；仅检查一个字符串也不能
区分初始化成功与发送路径中途失效。可观察协议应让不同执行阶段留下不同证据。
\end{failurebox}

\subsection{随机测试的边界}

随机测试不替代规格样例。固定种子保证失败可重现；生成器只改变复位位移并保持
其他布局合法；断言是“任何不指向 0xF000 的位移都必须被拒绝”。这种测试
检验的是性质，而不是堆积随机输入。

\subsection{完整验证命令}

\begin{lstlisting}[language=bash]
python3 tools/os.py verify
python3 tools/os.py test --layer unit
python3 tools/os.py test --layer integration
python3 tools/os.py test --layer randomized
python3 tools/os.py test --layer system
python3 tools/os.py test --layer failure-path
\end{lstlisting}

\subsection{裸机测试面对的特殊边界}

宿主应用测试失败时，运行库通常能打印异常和栈回溯。最早固件没有文件系统、
进程、调试服务甚至可用屏幕；三重故障还可能直接触发复位。因此测试设计必须
先创造观测通道，再利用外部进程限制时间并收集证据。

QEMU 进程退出码、串口字节、镜像哈希和反汇编都属于观测，但它们证明的范围
不同。串口出现 \code{RESET} 证明控制流至少到达相应发送点，不证明后面的
代码正确；镜像大小正确不证明复位跳转目标正确；单元测试通过也不证明 NASM
和链接脚本产生了预期字节。

\subsection{测试层次与故障定位}

\begin{longtable}{L{0.16\textwidth}L{0.26\textwidth}L{0.25\textwidth}L{0.19\textwidth}}
\toprule
\textbf{层次} & \textbf{主要对象} & \textbf{能够快速定位的故障} & \textbf{不能单独证明}\\
\midrule
\endfirsthead
\toprule
\textbf{层次} & \textbf{主要对象} & \textbf{能够快速定位的故障} & \textbf{不能单独证明}\\
\midrule
\endhead
单元测试 & 地址计算、解析和状态转换 & 边界、分支与错误分类 & 目标机器实际执行\\
集成测试 & 工具链组合和最终产物 & 节布局、符号、尺寸、固定字节 & 设备运行时行为\\
随机测试 & 满足生成约束的大量输入 & 手工样例遗漏的性质反例 & 未声明的全部输入空间\\
系统测试 & QEMU 中的完整镜像 & CPU、ROM、UART 的端到端协作 & 真实主板全部差异\\
失败路径 & 注入的设备或镜像异常 & 超时、拒绝和诊断可观测性 & 任意硬件故障\\
\bottomrule
\end{longtable}

测试层次越高，覆盖的组件越多，但失败定位越宽。低层测试提供快速而精确的
反馈，高层测试验证真实组合；二者不是替代关系。

\subsection{性质测试的生成域}

随机性只有与明确性质结合才有意义。复位向量测试先固定合法 ROM 尺寸、操作码
和其余字节，再从所有不能解码到 \code{0xF000} 的 16 位位移中抽样。若审计器
接受其中任意一个样本，就得到可保存的种子和最小失败输入。

固定随机种子保证 CI 与本地能够重现同一序列。随着阶段扩展，磁盘装载器可生成
截断 ELF、重叠段和溢出范围；页分配器可生成分配与释放操作序列；调度器可生成
任务状态转换。生成器必须维持输入域约束，否则测试只是在重复检查无意义垃圾。

\subsection{故障注入验证失败语义}

成功路径只验证设备按预期响应。系统工程还必须明确无响应、错误响应和部分完成
时的行为。当前失败镜像让发送过程持续观察不到 THRE，就绪轮询耗尽后返回失败。
它验证的是同一生产逻辑中的超时分支，而不是另写一个永远失败的测试替身。

后续 IDE 驱动应分别注入 BSY 不清除、ERR 置位和短读；内存分配器应覆盖耗尽
与重复释放；系统调用应覆盖非法地址和未知编号。每种故障都要有界终止，并留下
足以区分阶段的诊断信息。

\subsection{时间预算也是接口契约}

无限等待把硬件故障扩大成整个系统永久失去进展。轮询次数是目标端的逻辑预算，
QEMU 五秒总截止是宿主测试的外部保险；两者服务于不同层次。前者保证固件自行
返回失败，后者处理固件死循环、模拟器异常和测试框架错误。正常路径由最终
里程碑主动结束，不必机械耗尽整个截止时间。

时间阈值不能只按某台开发机的偶然速度选择。目标循环采用次数而非墙钟，系统
测试则给正常启动留下充足裕量。未来有定时器后，可把设备协议中的微秒或毫秒
上限转换为稳定的时钟刻度。

\subsection{证据链与回归保护}

每个里程碑保存从需求到测试的对应关系：需求定义可观察结果，设计说明状态
转换，代码实现机制，测试产生日志与产物，CI 在干净环境重放。回归发生时，
失败层次直接提示应先检查纯逻辑、链接布局还是整机交互。

\begin{keyidea}
测试数量不是目标。每个测试都应说明保护哪个不变量、在哪个边界观测，以及
失败时能够排除哪些原因。
\end{keyidea}
